\section{The Hermetic Method}

\label{sec:HermeticMethod}

\begin{quotex}
Ask, and it will be given you\newline
seek, and you will find;\newline
knock, and it will be opened to you. \flright{\itshape Luke 9:9} 
\end{quotex}

Hermetism differs from both religion and science, although it does not attempt to replace them. Unlike religion, which
accepts things on faith, and unlike science, which is based on facts and theories, Hermetism seeks knowledge. Not
knowledge of facts, no matter how mystical, sublime, or luminous, nor knowledge of articles of faith or scientific
knowledge, but knowledge of the mysteries, a direct, unmediated, intuitive knowledge.

Hermetism is a way or a path, so it must be traversed. Reading another book or article, or following another guru, will
not lead to Hermetic knowledge. Instead, spiritual practice is absolutely essential.

In the \textit{Letter IV on the Emperor}, \textbf{Valentin Tomberg} explains the essential prerequisite of spiritual exercises:

\begin{quotex}
One must establish the state of consciousness suitable for receiving revelations. 

\end{quotex}
This requires the following four steps:

\begin{enumerate}
\item The state of concentration without effort 
\item Watchful interior silence 
\item An inspired activity of imagination and thought where the conscious Self acts with the superconscious 
\item The conscious Self then stops its activity and contemplates everything which preceded 
\end{enumerate}
Tomberg recommends the \emph{Gospels}, \emph{John's Revelation} and parts of the Old Testament (e.g.
\emph{Genesis}, \emph{Ezekiel}) as suitable subjects for meditation or spiritual exercise, not only the Arcana of the
Tarot.

Now, it will happen that different Hermetists will develop skills in the respective areas they chose to meditate on. But
that is not what unites them, since only the depths of their meditations can do so. To further clarify this point,
Tomberg makes a distinction between the individual knowledge of Hermetists — which is knowledge of
particular facts — and their common knowledge, which is of the mysteries and involves depth of
consciousness.

For example, one may know his past lives as a fact, another may know about the chakras — that is their
individual interest and vocation. If I, for example, have no knowledge of past lives as a fact, then I simply ignore
it; it is not something I am obligated to believe. So Hermetism is not an accumulation of facts to be studied and
accepted uncritically.

Rather, Tomberg makes clear that Hermetism is a method. Just reading books will help no one become deep. Tomberg
explains:

\begin{quotex}
Spiritual exercises alone have taught [the initiate] — no theory or doctrine may in any way have rendered
him capable of knowing how to know. 

\end{quotex}
This will lead to a different kind of knowing, as represented by the Emperor. Hence,

\begin{quotex}
Hermetism does not teach what one ought to believe about man, God, and nature, but it teaches rather how to ask, seek
and knock in order to arrive at mystical experience, gnostic illumination, and the magical effect.

\end{quotex}

Only after having asked, sought, and knocked, can one know.

\flright{\itshape Posted on 2010-08-09 by Cologero}
